\documentclass[10pt,a4paper]{amsart}
\usepackage{amsmath,amssymb,amsthm,graphicx}
% Retained minimal preamble for the Stacks Project chapter more-algebra.tex (More on Algebra).
% The chapter file pinned-original-source.tex carries only \input{preamble} on its first line, so no
% preamble text lives inside the chapter itself.  The Stacks Project uses one single shared preamble for
% every chapter; the snapshot stages it identically as sources/stacks-chow/preamble.tex and
% sources/stacks-etale-cohomology/preamble.tex.  The block below is copied mechanically, line for line,
% from sources/stacks-chow/preamble.tex lines 14-18 (the Xy-pic setup) and 154-217 (the native theorem
% declarations together with the mathematical macros).  Nothing was added, removed or redefined; the
% wrapper class and the amsmath/amssymb/amsthm/graphicx base packages are supplied by the assembler.
% Chapter: more-algebra.tex -- More on Algebra.
% For commutative diagrams we use Xy-pic
\usepackage[all]{xy}

% We use 2cell for 2-commutative diagrams.
\xyoption{2cell}

% Theorem environments.
%
\theoremstyle{plain}
\newtheorem{theorem}[subsection]{Theorem}
\newtheorem{proposition}[subsection]{Proposition}
\newtheorem{lemma}[subsection]{Lemma}

\theoremstyle{definition}
\newtheorem{definition}[subsection]{Definition}
\newtheorem{example}[subsection]{Example}
\newtheorem{exercise}[subsection]{Exercise}
\newtheorem{situation}[subsection]{Situation}

\theoremstyle{remark}
\newtheorem{remark}[subsection]{Remark}
\newtheorem{remarks}[subsection]{Remarks}

\numberwithin{equation}{subsection}

% Macros
%
\def\lim{\mathop{\mathrm{lim}}\nolimits}
\def\colim{\mathop{\mathrm{colim}}\nolimits}
\def\Spec{\mathop{\mathrm{Spec}}}
\def\Hom{\mathop{\mathrm{Hom}}\nolimits}
\def\Ext{\mathop{\mathrm{Ext}}\nolimits}
\def\SheafHom{\mathop{\mathcal{H}\!\mathit{om}}\nolimits}
\def\SheafExt{\mathop{\mathcal{E}\!\mathit{xt}}\nolimits}
\def\Sch{\mathit{Sch}}
\def\Mor{\mathop{\mathrm{Mor}}\nolimits}
\def\Ob{\mathop{\mathrm{Ob}}\nolimits}
\def\Sh{\mathop{\mathit{Sh}}\nolimits}
\def\NL{\mathop{N\!L}\nolimits}
\def\CH{\mathop{\mathrm{CH}}\nolimits}
\def\proetale{{pro\text{-}\acute{e}tale}}
\def\etale{{\acute{e}tale}}
\def\QCoh{\mathit{QCoh}}
\def\Ker{\mathop{\mathrm{Ker}}}
\def\Im{\mathop{\mathrm{Im}}}
\def\Coker{\mathop{\mathrm{Coker}}}
\def\Coim{\mathop{\mathrm{Coim}}}

% Boxtimes
%
\DeclareMathSymbol{\boxtimes}{\mathbin}{AMSa}{"02}

%
% Macros for moduli stacks/spaces
%
\def\QCohstack{\mathcal{QC}\!\mathit{oh}}
\def\Cohstack{\mathcal{C}\!\mathit{oh}}
\def\Spacesstack{\mathcal{S}\!\mathit{paces}}
\def\Quotfunctor{\mathrm{Quot}}
\def\Hilbfunctor{\mathrm{Hilb}}
\def\Curvesstack{\mathcal{C}\!\mathit{urves}}
\def\Polarizedstack{\mathcal{P}\!\mathit{olarized}}
\def\Complexesstack{\mathcal{C}\!\mathit{omplexes}}
% \Pic is the operator that assigns to X its picard group, usage \Pic(X)
% \Picardstack_{X/B} denotes the Picard stack of X over B
% \Picardfunctor_{X/B} denotes the Picard functor of X over B
\def\Pic{\mathop{\mathrm{Pic}}\nolimits}
\def\Picardstack{\mathcal{P}\!\mathit{ic}}
\def\Picardfunctor{\mathrm{Pic}}
\def\Deformationcategory{\mathcal{D}\!\mathit{ef}}
\title{Krasner's lemma for a complete local domain of dimension one: a simple root survives every sufficiently small perturbation of the polynomial}
\author{The Stacks Project Authors}
\date{stacks-project:more-algebra.tex:local-snapshot-20261012; GNU FDL 1.2 or later}
\begin{document}
\maketitle
\setcounter{subsection}{1}
\noindent\textit{Editorial note.} All delivered mathematical text is the author's own, copied from the pinned Stacks Project chapter file more-algebra.tex; no statement or proof step was rewritten, reordered or removed. The optional title [Krasner's lemma] printed by the author on the statement environment is part of the counted statement span and is kept.\par\medskip
\noindent\textit{Editorial note.} The chapter file pinned-original-source.tex carries only an input of the shared preamble on its first line, so no preamble text lives inside the chapter. The retained file sources/stacks-more-algebra/retained-preamble.tex is a mechanical, line-for-line copy of the project's single shared preamble as staged in the snapshot, namely sources/stacks-chow/preamble.tex lines 14-18 (Xy-pic setup) and 154-217 (theorem declarations and mathematical macros). Nothing was added, removed or redefined, no package was added, and the excerpt serves only as preamble, never as a span coordinate.\par\medskip
\noindent\textit{Editorial note.} The source URL is the source-level Stacks Project URL, as required when more than one card is produced from one source; the per-claim tag is carried by claim\_key. The counted statement and proof contain no cross reference and no citation at all, so this card contributes no external reference stub and needs no bibliography; the section lead-in sentence that precedes the statement at source lines 33744-33745 is left outside the counted spans because it only announces the lemma.\par\medskip
\noindent\textit{Editorial note.} The proof writes the inverse of P'(alpha) as pi\textasciicircum{}\{-m\} a and uses the macro lim, which the retained preamble supplies; nothing outside the counted spans is required to read the delivered argument.\par\medskip
\noindent\textit{Allowed standard prerequisites (editorial).}\begin{itemize}
\item In a local domain of dimension one the maximal ideal is the radical of a principal ideal generated by any nonzero element of it. used at proof lines 33758-33759, where a nonzero pi in the maximal ideal is fixed and m = sqrt((pi)) is recorded. Turns the dimension-one hypothesis into the valuation-like statement that an element of the fraction field with a pole can be written as pi to a negative power times an element of A, which is what produces the integer m with P'(alpha)\textasciicircum{}\{-1\} = pi\textasciicircum{}\{-m\} a.
\item The Taylor expansion of a polynomial in two variables, P(x + y) = P(x) + P'(x) y + R(x, y) y\textasciicircum{}2 for some R in A[x, y]. recorded at proof lines 33763-33766 and used again at proof lines 33792-33796. Supplies the quadratic error term that makes the Newton correction work modulo the next power of the maximal ideal.
\item Completeness of a local ring with respect to its maximal-adic filtration: a sequence whose consecutive differences lie in higher and higher powers of the maximal ideal converges, and its limit can be used in A. used at proof line 33774, where beta is defined as the limit of the constructed sequence. Converts the inductively constructed sequence into an actual element beta of A with the required congruence beta - alpha in m\_A\textasciicircum{}c.
\end{itemize}\medskip
\begin{lemma}[Krasner's lemma]
\label{lemma-krasner}
Let $A$ be a complete local domain of dimension $1$. Let $P(t) \in A[t]$
be a polynomial with coefficients in $A$. Let $\alpha \in A$ be a root
of $P$ but not a root of the derivative $P' = \text{d}P/\text{d}t$.
For every $c \geq 0$ there exists an integer $n$ such that for any
$Q \in A[t]$ whose coefficients are in $\mathfrak m_A^n$ the polynomial
$P + Q$ has a root $\beta \in A$ with $\beta - \alpha \in \mathfrak m_A^c$.
\end{lemma}

\begin{proof}
Choose a nonzero $\pi \in \mathfrak m$. Since the dimension of $A$ is $1$
we have $\mathfrak m = \sqrt{(\pi)}$. By assumption we may write
$P'(\alpha)^{-1} = \pi^{-m} a$ for some $m \geq 0$ and $a \in A$.
We may and do assume that $c \geq m + 1$.
Pick $n$ such that $\mathfrak m_A^n \subset (\pi^{c + m})$.
Pick any $Q$ as in the statement. For later use we observe that we can write
$$
P(x + y) = P(x) + P'(x)y + R(x, y)y^2
$$
for some $R(x, y) \in A[x, y]$. We will show by induction that we can find a
sequence $\alpha_m, \alpha_{m + 1}, \alpha_{m + 2}, \ldots$ such that
\begin{enumerate}
\item $\alpha_k \equiv \alpha \bmod \pi^c$,
\item $\alpha_{k + 1} - \alpha_k \in (\pi^k)$, and
\item $(P + Q)(\alpha_k) \in (\pi^{m + k})$.
\end{enumerate}
Setting $\beta = \lim \alpha_k$ will finish the proof.

\medskip\noindent
Base case. Since the coefficients of $Q$ are in $(\pi^{c + m})$ we have
$(P + Q)(\alpha) \in (\pi^{c + m})$. Hence $\alpha_m = \alpha$ works.
This choice guarantees that $\alpha_k \equiv \alpha \bmod \pi^c$
for all $k \geq m$.

\medskip\noindent
Induction step. Given $\alpha_k$ we write
$\alpha_{k + 1} = \alpha_k + \delta$ for some $\delta \in (\pi^k)$.
Then we have
$$
(P + Q)(\alpha_{k + 1}) =
P(\alpha_k + \delta) + Q(\alpha_k + \delta)
$$
Because the coefficients of $Q$ are in $(\pi^{c + m})$ we see that
$Q(\alpha_k + \delta) \equiv Q(\alpha_k) \bmod \pi^{c + m + k}$.
On the other hand we have
$$
P(\alpha_k + \delta) =
P(\alpha_k) + P'(\alpha_k)\delta + R(\alpha_k, \delta)\delta^2
$$
Note that $P'(\alpha_k) \equiv P'(\alpha) \bmod (\pi^{m + 1})$
as $\alpha_k \equiv \alpha \bmod \pi^{m + 1}$. Hence we obtain
$$
P(\alpha_k + \delta) \equiv P(\alpha_k) + P'(\alpha) \delta
\bmod \pi^{k + m + 1}
$$
Recombining the two terms we see that
$$
(P + Q)(\alpha_{k + 1}) \equiv (P + Q)(\alpha_k) + P'(\alpha) \delta
\bmod \pi^{k + m + 1}
$$
Thus a solution is to take
$\delta = -P'(\alpha)^{-1} (P + Q)(\alpha_k) =  - \pi^{-m} a (P + Q)(\alpha_k)$
which is contained in $(\pi^k)$ by induction assumption.
\end{proof}

\end{document}
